% problem-000241 2 铝配合物结构与范德华力
\section*{2. 铝配合物结构与范德华力}
三(8-羟基喹啉)合铝(II)（AIO₃）配合物的化学式为 $\mathbf { C } _ { 2 7 } \mathbf { H } _ { 1 8 } \mathbf { N } _ { 3 } \mathbf { O } _ { 3 } \mathbf { A } \mathbf { l }$ ，是一种优异的小分子半导体材料，可用于有机显示等领域。

\subsection*{2-1}
$\mathsf { A l Q } _ { 3 }$ 分子有两种典型异构体，如图2.1所示。写出它们所属的分子点群。

\subsection*{2-2}
预测 $\mathsf { A l Q } _ { 3 }$ 配合物中 $\mathsf { A l } - \mathsf { O }$ 与 $\mathsf { A l - N }$ 键长的大小关系，解释原因。

$2 - 3 \mathsf { A l Q } 3$ 分子可通过范德华力形成分子晶体。一般地，以范德华力主导的分子间相互作用势能E(l)可用 Lennard-Jones 关系表示

$
E (l) = \frac {a}{l ^ {1 2}} - \frac {b}{l ^ {6}}
$

其中，l是分子质心之间的距离，a和b均为大于0的常数。画出E(l)-l的关系图，求 $\mathsf { A l Q } _ { 3 }$ 分子的范德华半径。

\subsection*{2-4}
$\mathsf { A l Q } _ { 3 }$ 分子可结晶为三斜晶体，一个晶胞里有两个不等价的 $\mathsf { A l Q } _ { 3 }$ 分子。其晶胞参数为 $a = 6 . 1$ 181 Å, b = 13.27 A, c= 14.43 Å, $a = 6 6 . 0 6 ^ { \circ }$ $\beta = 8 8 . 5 6 ^ { \circ }$ $\gamma = 8 4 . 0 3 ^ { \circ }$ 。计算 $\mathsf { A l Q } _ { 3 }$ 分子晶体的密度（单位： $\mathbf { g } \mathsf { c m } ^ { - 3 } )$

\subsection*{2-5}
研究发现，对8-羟基喹啉配体在羟基对位进行修饰，所得 $A l Q _ { 1 }$ 衍生物分子的发射光谱可在整个可见光范围内调制。对下列衍生物分子发射光谱的峰值波长排序。

（图：）
